\source{Лекции, с. 30–31; заключительный подсчёт восстановлен составителем; Карпов, §3.1}
\begin{definition}[Гамильтонов граф]
Гамильтонов путь — простой путь, содержащий все вершины графа.
Гамильтонов цикл — простой цикл, содержащий все вершины.
Граф называется гамильтоновым, если имеет гамильтонов цикл; далее \(n\ge3\).
Степенная последовательность — степени в неубывающем порядке
\(d_1\le\cdots\le d_n\).
\end{definition}
\begin{theorem}[Хватал]
Пусть \(n\ge3\), \(d_1\le\cdots\le d_n\). Если для каждого целого
\(1\le i<n/2\) выполнено
\[
\boxed{d_i\le i\quad\Longrightarrow\quad d_{n-i}\ge n-i,}
\]
то граф имеет гамильтонов цикл.
\end{theorem}
\begin{proof}
Предположим противное. Добавление рёбер не уменьшает ни одну
из упорядоченных степеней, поэтому сохраняет условие: если новая
\(d_i\le i\), то и старая \(d_i\le i\), а новая \(d_{n-i}\) не меньше старой.
Добавим к исходному графу как можно больше рёбер, сохраняя негамильтоновость;
получим \(H\). Он не полный, а добавление любого отсутствующего ребра
делает его гамильтоновым. В гамильтоновом цикле \(H+uv\) обязательно
используется новое ребро, иначе \(H\) уже гамильтонов.
Удалив его, получаем гамильтонов путь
\(x_1=u,x_2,\ldots,x_n=v\) в \(H\).

Среди несмежных \(u,v\) выберем пару с максимальной суммой степеней,
\(h=d_H(u)\le d_H(v)=q\). Рассмотрим
\[
S=\{i:x_i v\in E(H)\},\qquad
T=\{i:u x_{i+1}\in E(H)\}\subseteq\{1,\ldots,n-1\}.
\]
Имеем \(|S|=q\), \(|T|=h\). Если \(i\in S\cap T\), возникает цикл
\[
u,x_2,\ldots,x_i,v,x_{n-1},\ldots,x_{i+1},u,
\]
проходящий через все вершины. Крайние вырожденные индексы исключены
несмежностью \(u,v\). Следовательно, \(S\cap T=\varnothing\), и
\(h+q\le n-1\), в частности \(h<n/2\).

Каждый несосед \(v\), отличный от \(v\), имеет степень не больше \(h\),
иначе выбранная сумма не максимальна. Таких вершин \(n-1-q\ge h\).
Каждый несосед \(u\) имеет степень не больше \(q\le n-h-1\).
Вместе с \(u\) это \(n-h\) вершин степени не больше \(n-h-1\)
(\(h\le n-h-1\)). Если \(h\ge1\), то
\[
d_h(H)\le h,\qquad d_{n-h}(H)\le n-h-1,
\]
что противоречит условию Хватала. Если \(h=0\), вершина \(u\) изолирована,
а все остальные степени не больше \(n-2\). Тогда
\(d_1(H)=0\le1\), \(d_{n-1}(H)\le n-2\), снова противоречие при \(i=1<n/2\).
\end{proof}
\begin{center}
\begin{tikzpicture}[>=Stealth,scale=.9,every node/.style={circle,fill=cbBlue,inner sep=2pt}]
\foreach \x/\name in {0/u,1.1/a,2.2/b,3.3/c,4.4/d,5.5/v}{\node (\name) at (\x,0) {};}
\draw (u)--(a)--(b)--(c)--(d)--(v);
\draw[cbOrange,thick] (u) to[bend left=35] (c);
\draw[cbGreen,thick] (b) to[bend right=35] (v);
\node[fill=none,below] at (u) {$u=x_1$};\node[fill=none,below] at (b) {$x_i$};
\node[fill=none,below] at (c) {$x_{i+1}$};\node[fill=none,below] at (v) {$v=x_n$};
\end{tikzpicture}
\end{center}
Два выделенных ребра при \(i\in S\cap T\) замыкают путь в гамильтонов цикл.

\textbf{Книжное дополнение для следующего билета.}
\begin{lemma}[Добавление ребра]\label{lem:closure}
Если \(u,v\) несмежны и \(d(u)+d(v)\ge n\), то \(G\) гамильтонов
тогда и только тогда, когда \(G+uv\) гамильтонов.
\end{lemma}
\begin{proof}
Только обратное направление требует проверки. Если цикл использует
\(uv\), удаляем его и получаем гамильтонов путь. Для множеств \(S,T\)
выше сумма мощностей не меньше \(n\), поэтому они пересекаются в
множестве из \(n-1\) индексов. Та же перестройка даёт цикл в \(G\).
\end{proof}
\begin{definition}[Замыкание; дополнение по книге]
Последовательно добавляя рёбра между несмежными вершинами с суммой
степеней не меньше \(n\), получаем замыкание \(\cl(G)\).
Оно одновременно гамильтоново с \(G\) и не зависит от порядка добавлений.
\end{definition}
\begin{proof}[Обоснование]
Первое следует из леммы. Для двух завершённых процессов \(H_1,H_2\)
индукцией по добавлениям первого покажем \(H_1\subseteq H_2\).
Если предыдущие рёбра уже есть в \(H_2\), степени концов следующего
там не меньше, чем при его добавлении. Отсутствие ребра в \(H_2\)
противоречило бы завершённости второго процесса. Симметрично получаем
обратное включение. В частности, \(\cl(G)=K_n\) влечёт гамильтоновость.
\end{proof}
\textbf{Дополнительное свойство степенного критерия (по книге).}
\begin{cor}[Точность критерия как условия на степени]
Если неубывающая последовательность \(0\le a_1\le\cdots\le a_n\le n-1\)
нарушает условие Хватала при \(i=h\), существует негамильтонов граф,
степени которого покоординатно не меньше \(a_1,\ldots,a_n\).
\end{cor}
\begin{proof}
Возьмём три множества \(A,B,D\) размеров \(h,h,n-2h\).
В \(B\cup D\) проведём все рёбра; соединим каждую вершину \(A\)
со всеми вершинами \(B\), других рёбер у \(A\) нет.
Степени равны \(h\) на \(A\), \(n-h-1\) на \(D\), \(n-1\) на \(B\),
поэтому они мажорируют заданную последовательность.
После удаления \(B\) остаются \(h+1\) компонент: отдельные вершины \(A\)
и непустая клика \(D\). Удаление \(h\) вершин из гамильтонова цикла
оставляет не более \(h\) непустых отрезков, значит, в гамильтоновом графе
компонент после такого удаления не больше \(h\). Противоречие.
\end{proof}
Критерий достаточен для конкретного графа, но не необходим: например,
цикл \(C_5\) гамильтонов, а при \(i=2\) его степени нарушают условие.
